
\section{Finite Sample Guarantee of MSQP}\label{sec:theory}


In this section, we present our primary sample complexity guarantee for our  MSQP algorithm under the assumptions of realizability of  $q^{\star}_{\alpha}$ and $l^{\star}_{\alpha}$  and partial coverage. We first show the learned policy $\hat \pi_{\alpha}$ can compete with $\pi^{\star}_{\alpha}$.  Finally we show  $\hat \pi_{\alpha}$ can compete with $\pi^{\star}$ by selecting $\alpha$ properly.   

We first introduce the flattened behavior policy $\pi^{\diamond}_b$, which is uniform on the support of $\pi_b$.
We use it as a technical device to define a model-free concentrability coefficient following \citet{xie2021bellman}. 
\begin{definition}[Model-free concentrability coefficient]
Define
\begin{align*}
    C_{\Qcal,d_{\pi^{\star}_{\alpha},\mu_0} }\coloneqq \sup_{q\in \Qcal}\frac{\EE_{s \sim d_{\pi^{\star}_{\alpha},\mu_0},a\sim \pi^{\diamond}_b(a\mid s) }[\|q(s,a) -q^{\star}_{\alpha}(s,a)\|^2_2 ]}{\EE_{(s,a)\sim P_b}[\|q(s,a) -q^{\star}_{\alpha}(s,a)\|^2_2]}
\end{align*}
where $\pi^{\diamond}_b(\cdot \mid s)=\begin{cases}
     0  \quad & \pi_b(\cdot \mid s) = 0 \\
     1/|\{a \in \Acal\mid \pi_b(a \mid s)>0\}|  \; & \pi_b(\cdot \mid s)>0
    \end{cases}$ is the flattened behavior policy.
\end{definition}
Clearly, $C_{\Qcal,d_{\pi^{\star}_{\alpha},\mu_0} }$ is smaller than density-ratio-based concentrability coefficient, in other words, 
$$C_{\Qcal,d_{\pi^{\star}_{\alpha},\mu_0} }\leq \max_{(s,a)} \frac{d_{\pi^{\star}_{\alpha},\mu_0}(s)\pi^{\diamond}_b(a\mid s)}{P_b(s)\pi_b(a\mid s)}.
$$ 
Here, we always have $\|\pi^{\diamond}_b/\pi_b\|< \infty$ even if $\pi_b(a\mid s)$ is $0$ for some $(s,a)$. In the special case where $\pi_b(a\mid s)\geq 1/C'$ for any $(s,a)$, we have $C_{\Qcal,d_{\pi^{\star}_{\alpha},\mu_0} }\leq C'\|d_{\pi^{\star}_{\alpha},\mu_0}/P_b\|_{\infty} $. 
The coefficient $C_{\Qcal,d_{\pi^{\star}_{\alpha},\mu_0} }$ is is a refined concentrability coefficient, which adapts to a function class $\Qcal$. For example, in linear MDPs, it reduces to a relative condition number as follows. Similar properties are obtained in related works \citep{xie2021bellman,uehara2021pessimistic}. 

\begin{example}[Linear MDPs]
A linear MDP is one such that, for a known feature vector $\phi:\Scal \times \Acal \to \RR^d$, the true density satisfies $P(s' \mid s,a)=\langle \mu^{\star}(s'),\phi(s,a)\rangle $ for some $\mu^{\star}:\Scal \to \RR^d$ and the reward function satisfies $\EE[r\mid s,a]=\langle \theta_r, \phi(s,a) \rangle $ for some $\theta_r \in \RR^d$. 

In linear MDPs, $q^{\star}_{\alpha}$ is clearly linear in $\phi(s,a)$. Hence, the natural function class is $\Qcal=\{\langle \theta,\phi(s,a)\rangle \mid \|\theta\|\leq B\}$ for a certain $B\in \mathbb{R}^{+}$. Then, we have 
\begin{align*}
C_{\Qcal,d_{\pi^{\star}_{\alpha},\mu_0} }=\sup_{x\neq 0}\frac{x^{\top}\EE_{s \sim d_{\pi^{\star}_{\alpha},\mu_0}, a\sim \pi^{\diamond}_b(a\mid s) }[\phi(s,a)\phi(s,a)^{\top}]x  }{x^{\top}\EE_{(s,a)\sim P_b}[\phi(s,a)\phi(s,a)^{\top}]x}. 
\end{align*}
\end{example}


We are now prepared to present our main result, which states that given the realizability of the soft Q-function $q^{\star}_{\alpha}$ and Lagrange multiplier $l^{\star}_{\alpha}$, it is possible to compete with $\pi^{\star}_{\alpha}$ under the coverage condition $C_{\Qcal,d_{\pi^{\star}_{\alpha},\mu_0} }<\infty,\|d_{\pi^{\star}_{\alpha},P_b}/P_b\|_{\infty}<\infty $. 

\begin{theorem}[$\hat \pi_{\alpha}$ can compete with $\pi^{\star}_{\alpha}$]\label{thm:main}
Fix $\alpha>0$. Suppose Assumptions \pref{assum:scale2}, \pref{assum:scale}, \pref{assum:softmax}, and \pref{assum:largrange} hold. 
With probability $1-\delta$, the regret $ J(\pi^{\star}_{\alpha})-J(\hat \pi_{\alpha})$  is upper-bounded by 
\begin{align*}
n^{-1/4}\mathrm{Poly}\left(|\Acal|,\Bcal_{\Qcal},\Bcal_{\Lcal},C_{\Qcal,d_{\pi^{\star}_{\alpha},\mu_0} }, \ln\left(\frac{|\Qcal||\Lcal|}{\delta}\right),\frac{1}{1-\gamma},R_{\max}\right)
\end{align*}

\end{theorem}
The proof mainly consists of two steps: (1) obtaining $L^2$-errors of $\hat q_{\alpha}$ as previously demonstrated in \pref{thm:convergence_soft}, (2) translating this error into the error of $\hat \pi_{\alpha}$. In the second step, the Lipshitz continuity of the softmax function plays a crucial role. If there is no regularization ($\alpha=0$) and the greedy policy of $q^{\star}_{0}$ is utilized, the second step does not proceed (without any further additional assumptions). 

Our ultimate goal is to compete with $\pi^{\star}$. 
\pref{thm:main} serves as the primary foundation for this goal. The remaining task is to analyze the approximation error $J(\pi^{\star})-J(\pi^{\star}_{\alpha})$. Fortunately, this term can be controlled through  $\alpha$ and the density ratio between $\pi^{\star}$ and $\pi_b$. Then, by properly controlling $\alpha$, we can obtain the following sample complexity result.
 
\begin{theorem}[PAC guarantee of $\hat \pi_{\alpha}$]\label{cor:main}
Fix any $\epsilon>0$. Suppose Assumptions~\pref{assum:scale2}, \pref{assum:scale}, \pref{assum:softmax}, and \ref{assum:largrange} hold for $\alpha=c/n^{1/8}$ and  $\|\pi^{\star}_{0}/\pi_b\|_{\infty}\leq C_{0}$, $C_{\Qcal,d_{\pi^{\star}_{\alpha},\mu_0} }<\infty$. Then, if $n$ is at least  
{
\begin{align*}
\epsilon^{-8}\mathrm{Poly}(|\Acal|,\Bcal_{\Qcal},\Bcal_{\Lcal},C_{\Qcal,d_{\pi^{\star}_{\alpha},\mu_0} }, \ln(|\Qcal||\Lcal|/\delta),%
(1-\gamma)^{-1},\ln(C_0),R_{\max}), 
\end{align*}
}
with probability at least $1-\delta$, we can ensure $ J(\pi^{\star})-J(\hat \pi_{\alpha})\leq \epsilon$. 
\end{theorem}


{\finaledit In summary, the realiazability of $q^{\star}_{\alpha},l^{\star}_{\alpha}$, per-step coverage $\|\pi^{\star}_0/\pi_b\|_{\infty}<\infty$ and partial coverage $C_{\Qcal,d_{\pi^{\star}_{\alpha},\mu_0} }<\infty,\|d_{\pi^{\star}_{\alpha},P_b}/P_b\|_{\infty}<\infty $ are sufficient to compete with $\pi^{\star}$.} This is a novel and attractive result. Firstly, if we solely use the na\"ive FQI or Bellman residual minimization, existing PAC results require the global coverage $\|d_{\pi,\mu_0}/P_b\|_{\infty}<\infty$ for any possible policy $\pi$ \citep{munos2008finite,antos2008learning}. Our result only requires coverage under a single policy $\pi^{\star}_{\alpha}$ (near-optimal policy). Secondly, we only require the realizability of two functions, and we do not necessitate realizability-type conditions for all policies in the policy class or Bellman completeness, unlike existing works with partial coverage \citep{xie2021bellman,jiang2020minimax}. 



The most similar result is \citet{zhan2022offline}. However, our guarantee possesses a certain advantage over their guarantee as follows. They demonstrate the realizability of certain functions $\tilde w^{\star}_{\alpha},v^{\star}_{\alpha}$ and partial coverage $\|d_{\tilde\pi_{\alpha},\mu_0}/P_b\|< \infty$ are sufficient conditions in offline RL, where $\tilde w^{\star}_{\alpha}=d_{\tilde\pi_{\alpha},\mu_0}/P_b$ ($\tilde\pi_{\alpha}$ is a certain regularized optimal policy, but fundamentally distinct from $\pi^{\star}_{\alpha}$) and $v^{\star}_{\alpha}$ is a near-optimal regularized value function parameterized by $\alpha$. Here, we have  $\tilde w^{\star}_{0}=w^{\star},v^{\star}_{0}=v^{\star}$. Our guarantee has a similar flavor in the sense that it roughly illustrates realizability and partial coverage are sufficient conditions. 
However, the meanings of realizability and partial coverage are significantly different. In particular, by employing our algorithm, we can ensure PAC guarantees under the boundedness of the refined concentrability coefficient $C_{\Qcal,d_{\pi^{\star}_{\alpha},\mu_0}}<\infty$ (and $\|d_{\pi^{\star}_{\alpha},P_b}/P_b\|_{\infty}$ through $\Bcal_{\Lcal}$). 
As a result, the $L^{\infty}$-norm of the density-ratio-based concentrability coefficient $\|d_{\pi^{\star}_{\alpha},\mu_0}/P_b\|_{\infty}$ can even be infinite. More specifically, we can permit situations where $\max_{s}\mu_0(s)/P_b(s)=\infty$ as we will see the practical example soon. Conversely, \citet{zhan2022offline} excludes this possibility since the algorithm explicitly estimates the density ratio $\tilde w^{\star}_{\alpha}$. 


















\begin{example}[Contextual bandit under external validity]
 \label{exa:external}
 We consider the contextual bandit setting where we want to optimize $J(\pi)=\EE_{s\sim \mu_0,a\sim \pi(s),r\sim P_r(s,a)}[r]$ using offline data $s \sim P_b,a\sim \pi_b(s),r\sim P(s,a)$. This is the simplest RL setting with $\gamma=0$. Here, note $\mu_0$ could be different from $P_b$. This case often happens in practice as discussed in the literature on causal inference related to external validity \citep{pearl2022external,dahabreh2019generalizing,uehara2020off}, which refers to the shift between the target population and the offline data. Here, our PAC guarantee does not require that $\mu_0(s)$ is covered by $P_b(s)$ in terms of the density ratio as long as the relative condition number is upper-bounded when we use linear models. On the other hand, \citet{zhan2022offline} excludes this possibility.  
\end{example}

{\finaledit Despite the aforementioned advantage of our approach, unfortunately, our sample complexity of $O(1/\epsilon^8)$ is slower compared to that of $O(1/\epsilon^6)$ in \citet{zhan2022offline}. In the following, we demonstrate that MQP, which is a special version of MSQP when $\alpha \to 0$, can achieve a faster rate of $O(1/\epsilon^2)$.}


\section{Finite Sample Guarantee of MQP} \label{sec:margin}


In this section, building upon the convergence result of $\hat q_0$, we demonstrate the finite sample guarantee of our MQP algorithm under partial coverage. We first introduce the soft margin.

\begin{assum}[Soft margin]\label{assum:margin}
For any $a'\in \Acal$, there exists $t_0 \in \mathbb{R}^{+},\beta \in (0,\infty]$ such that 
\begin{align*}
    \PP_{s \sim d_{{\pi^{\star},\mu_0}}}(0<|q^{\star}(s,\pi^{\star}(s))-q^{\star}(s,a')|<t)\leq (t/t_0)^{\beta} 
\end{align*}
for any $t>0$. Here, we use the convention $x^{\infty}=0$ if $0<x<1$ and $x^{\infty}=\infty$ if $x>1$. 
\end{assum}

In the extreme case, if there exists a gap in $q^{\star}$ (also known as a hard margin) so that the best action is always better than the second-best by some lower bounded amount, then the soft margin is satisfied with $\beta=\infty$. Thus, the soft margin is more general than the gap condition  used in \citet{simchowitz2019non,wu2022gap}.
Crucially, a gap generally does \textit{not} exist in continuous state spaces unless Q-functions are discontinuous or one action is always option, or a gap involves a large $t_0$ constant in discrete state spaces with bad dependence on the number of states. In contrast, a soft margin with some $\beta>0$ generally holds (see, \eg, lemma 4 in \citealp{hu2021fast}). The soft margin is widely used in the literature on classification, decision making, and RL \citep{audibert2007fast,perchet2013multi,luedtke2020performance,hu2021fast,hu2022fast}.






\begin{theorem}[PAC guarantee of $\hat \pi_0$]\label{thm:pac_q_function}
Suppose Assumptions~\ref{assum:q_star}, \ref{assum:lagrange2}, and \ref{assum:margin} hold and  $\|\pi^{\star}/\pi_b\|_{\infty}\leq C_0$. 
 Fix any $\epsilon>0$. Then, if $n$ is at least 
\begin{align*}
\textstyle
\{\frac{|\Acal|}{\epsilon}\}^{\frac{4+2\beta}{\beta}} \mathrm{Poly}\left(t^{-1}_0,|\Acal|,\Bcal_{\Qcal},\Bcal_{\Lcal},C_{\Qcal,d_{\pi^{\star},\mu_0} }, \ln\left (\frac{|\Qcal||\Lcal|}{\delta}\right) %
,(1-\gamma)^{-1},\ln(C_0),R_{\max} \right)%
\end{align*}
with probability at least $1-\delta$, we can ensure $ J(\pi^{\star})-J(\hat \pi_{0})\leq \epsilon$. 
\end{theorem}

The proof mainly consists of two steps: (1) obtaining $L^2$-errors of $\hat q_{0}$ as demonstrated in \pref{thm:convergence_q}, (2) translating this error into the error of $\hat \pi_{0}$. In the second step, the soft margin plays a crucial role.

These theorems indicate that the realizability of the $Q$-function $q^{\star}$ and Lagrange multiplier $l^{\star}$, and the soft margin are sufficient for the PAC guarantee under partial coverage $C_{\Qcal,d_{\pi^{\star},\mu_0} }<\infty$ {\finaledit , $\|d_{\pi^{\star},P_b}/P_b\|_{\infty}<\infty$ }. Our algorithm is \emph{agnostic} to $\beta$ and operates under any value of $\beta$. {\finaledit In particular, when there is a gap ($\beta=\infty$), we can achieve sample complexity of $O(1/\epsilon^2)$ \footnote{{Similar to the findings in \citet{wang2022gap}, in general offline RL, we may potentially achieve a result of $O(1/\epsilon)$. We leave room for further enhancements in future research. }
}.} In comparison to \pref{thm:main}, although we additionally use the soft margin, the realizability in \pref{thm:pac_q_function} is more appealing since it is imposed on the standard Q-function $q^{\star}$. 
The closest guarantee to our work can be found in \citet{chen2022offline}, which demonstrates that the existence of the gap in $q^{\star}$, the realizability of $q^{\star},w^{\star}(:=d_{\pi^{\star},\mu_0}/P_{\pi_b})$, and partial coverage $\|w^{\star}\|_{\infty}< \infty$ are sufficient conditions. { A similar comparison is made in \citet{ozdaglar2023revisiting}.} In comparison to their work, we use the soft margin, which is significantly less stringent. 







